\section{Jordan's theorem supplies the standard-kernel hypothesis}
\label{sec:standard-jordan-kernel}

\paragraph{Status and provenance.}
HD-SS07 is a separate strengthening within reserved question Q2ed.
The kernel dichotomy and the application of the general Jordan theorem
were proposed by \texttt{/root/g34\_eightfolds}; root independently
considered the same bound and supplied the primitive-only alternative
recorded below. This is not an application of a primitive-group bound
to a reducible representation. The complete source-access and
proof-reading scopes are recorded separately. Complete author, root and independent whole-proof readings pass; the exact source and computation scopes are retained in their records.

\begin{lemma}[HD-SS07]
Let \(J\) be finite, and suppose
\[
 V=(\operatorname{Std}_n\circ x)\otimes U
 \quad\hbox{is faithful},\qquad
 x:J\longrightarrow S_n,\quad A_n\subset x(J),\quad
 n\ge5,\quad b=\dim U\ge2,
 \tag{SSJ.1}
\]
where \(U\) is an actual finite linear representation.
Then exactly one of the following holds:
\[
 \ker U\hbox{ is transitive on the letters},\qquad
 \ker U=1\ \hbox{ and }\ n\le\max(b+1,72).
 \tag{SSJ.2}
\]
Consequently in the range \(n>168(b+1)\) the kernel-transitivity
hypothesis of HD-SS06 is automatic. All conclusions of that theorem
then hold under its remaining exact-\(K\), multiplicity and
period-rank hypotheses, without an added primitivity assumption.

In particular, a source of minimal coefficient degree among actual
exact-\(K\), rank-nine naturally polarized sources of a fixed signed
determinant type cannot have a stable standard factorization with
\[
                         b\ge2,\qquad n>168(b+1).
 \tag{SSJ.3}
\]
The restriction in (SSJ.3) concerns such a minimal source. It is
not an exclusion of all nonminimal sources, and the descent does
not become an irreducible induction.
\end{lemma}

\begin{proof}
Set \(K_0=\ker U\), \(P=x(J)\), and \(R=x(K_0)\lhd P\).
As \(A_n\) is nonabelian simple, \(R\cap A_n\) is either all
of \(A_n\) or trivial. In the latter case
\([R,A_n]\subset R\cap A_n=1\). The centralizer of \(A_n\)
in \(S_n\) is trivial, so \(R=1\). To see that centralizer
fact directly, a permutation commuting with the transitive
\(A_n\) is determined by one image; the stabilizer \(A_{n-1}\)
fixes only that letter for \(n\ge5\), forcing that image to be
the same letter.
If \(R=1\), then \(K_0\) acts trivially on both tensor factors,
and faithfulness of \(V\) gives \(K_0=1\). Otherwise \(R\)
contains \(A_n\) and is transitive. The two alternatives are
disjoint.

Suppose now that \(K_0=1\), so \(U\) embeds \(J\) faithfully
in \(\operatorname{GL}_b(\mathbf C)\). Put
\[
                  d=\max(b,71),\qquad
                  \widetilde U=U\oplus1^{\oplus(d-b)}.
 \tag{SSJ.4}
\]
This is a faithful embedding in dimension \(d\). It need not be
irreducible or primitive. Collins' \emph{general} Theorem~A
\cite[p.~412]{HD:Collins2007} applies to every finite subgroup
of \(\operatorname{GL}_d(\mathbf C)\), \(d\ge71\), and supplies
an abelian normal subgroup \(A\lhd J\) such that
\[
                            [J:A]\le(d+1)!.
 \tag{SSJ.5}
\]
The theorem's complete reduction and proof explicitly include
reducible and imprimitive groups: Section~2, Lemma~1,
Theorem~2 and Corollary~3, pp.~414--417, and Sections~3--4,
pp.~417--423. We use no assertion that \(\widetilde U\)
is primitive.

The image \(x(A)\) is abelian and normal in \(P\). Its intersection
with \(A_n\) is trivial, and the same centralizer argument as above
gives \(x(A)=1\). Thus \(J/A\) surjects onto \(P\), and
\[
                    \frac{n!}{2}\le |P|\le[J:A]\le(d+1)!.
 \tag{SSJ.6}
\]
If \(n\ge d+2\), its left side is at least
\((d+2)!/2>(d+1)!\), a contradiction. Hence
\(n\le d+1=\max(b+1,72)\).

For \(b\ge2\), \(168(b+1)\) exceeds both \(b+1\) and \(72\).
Therefore the range in (SSJ.3) forces the first alternative of
(SSJ.2). HD-SS06 supplies an actual same-field, same-determinant,
rank-nine source of strictly smaller coefficient degree. This
proves the minimal-source assertion with precisely its stated
comparison class.
\end{proof}

\paragraph{A sharper optional bound when the partner is primitive.}
There is an independent way to check the main primitive-source
application without general Jordan theory. If \(V\) is primitive,
then \(U\) is primitive: a direct-sum system of imprimitivity for
\(U\) would give one for \(V\) after tensoring each summand with
\(\operatorname{Std}_n\).
When \(U\) is faithful and primitive, Collins' published primitive
Theorem~A \cite[p.~761]{LQ:Collins2008} gives
\([J:Z(J)]\le(b+1)!\) for \(b\ge13\).
Since \(x(Z(J))=1\), this implies \(n\le b+1\).
For \(2\le b\le12\), the largest listed exceptional bound is
\(448345497600\), and
\[
                448345497600<15!/2=653837184000.
 \tag{SSJ.7}
\]
The nonexceptional degrees ten and eleven have smaller
\((b+1)!\) bounds. Thus \(n\le14\) in this range and
\(n\le\max(b+1,14)\) throughout the primitive case.
This sharper statement is optional: the general proof above uses
only the 2007 theorem and permits a nonprimitive partner.

\paragraph{Exact source-access boundary.}
The 2007 theorem and its complete supplied proof were read in a
transcription of the actual published 13-page paper; the publisher
issue page and Oxford records independently confirm the authorship,
publication and general \(d\ge71\) statement. The publisher PDF
redirected to an access/verification page, and the Oxford record
offered metadata rather than a public file. No binary-PDF retrieval
or visual inspection of that 2007 paper is claimed. The full
transcription and original web responses are retained locally.
